% ============================================================================
% CAPÍTULO 5: MATERIALES Y MÉTODOS
% ============================================================================
%
% REQUISITOS SEGÚN LA FACULTAD:
%
% DISEÑO DEL ESTUDIO:
% • Tipo de investigación, enfoque metodológico
% • Explicar de manera detallada el diseño del estudio acorde a la pregunta y
%   objetivos de investigación con sustento teórico
%
% POBLACIÓN DE ESTUDIO:
% • Descripción de la población
% • Criterios de selección
% • Tamaño de la muestra y tipo de muestreo
%
% DEFINICIÓN OPERACIONAL DE VARIABLES:
% • Se presentan las variables del estudio
% • Tabla con definición semántica de cada variable (dependiente, independiente,
%   universales, confusores, predictores, intervinientes)
% • Definición operacional
% • Dimensiones
% • Escalas de medición apropiadas para cada tipo de variable
%
% TÉCNICAS Y PROCEDIMIENTO:
% • Describir técnicas e instrumentos empleados
% • Análisis de contenido, confiabilidad, consistencia interna y validación de
%   constructo (si aplica muestra amplia)
% • Expresar claramente procedimientos de reclutamiento, medición de variables,
%   técnicas analíticas, bioquímicas, clínicas u otras
% • En estudios de intervención: indicar asignación de grupos
% • Descripción del análisis de información por niveles/etapas/fases (univariado,
%   bivariado, multivariado)
%
% CONSIDERACIONES ÉTICAS:
% • Presentar aspectos éticos según principios básicos y lineamientos internacionales,
%   nacionales y locales
% • Carta de consentimiento informado en anexos
%
% NOTA IMPORTANTE: La redacción debe estar en TIEMPO PASADO (ya se realizó)
% ============================================================================

\chapter{Materiales y Métodos}
\label{cap:metodologia}

% ----------------------------------------------------------------------------
% SECCIÓN 5.1: Diseño del Estudio
% ----------------------------------------------------------------------------
\section{Diseño del Estudio}
\label{sec:diseño_estudio}

% INSTRUCCIONES:
% - Especifica: tipo de investigación, enfoque, diseño específico
% - Explica DETALLADAMENTE con sustento teórico
% - Justifica por qué este diseño responde a tu pregunta de investigación

\subsection{Tipo de Investigación y Enfoque}
\label{subsec:tipo_investigacion}

Este estudio se desarrolló bajo un enfoque \textit{[cuantitativo/cualitativo/mixto]} con un tipo de investigación \textit{[descriptivo/correlacional/explicativo/exploratorio]}.

\textbf{Justificación del enfoque:}

El enfoque \textit{[tipo]} fue seleccionado porque \textit{[justificación basada en tu pregunta y objetivos]} \cite{autor_metodologia}.

\subsection{Diseño Específico del Estudio}
\label{subsec:diseño_especifico}

El diseño de investigación fue \textit{[especifica el diseño]}:

\textbf{Ejemplos según tipo:}
\begin{itemize}
    \item \textbf{Experimental:} estudio experimental aleatorizado/cuasi-experimental
    \item \textbf{Observacional:} transversal/longitudinal (cohorte/casos y controles)
    \item \textbf{Descriptivo:} encuesta transversal/estudio de casos
\end{itemize}

\textbf{Características del diseño:}

\begin{itemize}
    \item \textbf{Temporalidad:} \textit{[Transversal/Longitudinal/Prospectivo/Retrospectivo]}
    \item \textbf{Intervención:} \textit{[Observacional/Experimental]}
    \item \textbf{Grupos de comparación:} \textit{[Ninguno/Grupo control/Múltiples grupos]}
    \item \textbf{Momento de medición:} \textit{[Una sola vez/Múltiples mediciones]}
\end{itemize}

\textbf{Esquema del diseño:}

[Si es posible, incluye un diagrama del diseño]

% \begin{figure}[htbp]
%     \centering
%     \includegraphics[width=0.7\textwidth]{figuras/esquema_diseño.png}
%     \caption{Esquema del diseño de investigación}
%     \label{fig:esquema_diseño}
% \end{figure}

% INSTRUCCIÓN: Crea tu diagrama de diseño y colócalo en figuras/esquema_diseño.png
% Luego descomenta las líneas de arriba

\subsection{Justificación del Diseño}
\label{subsec:justificacion_diseño}

Este diseño fue seleccionado porque \textit{[explicación detallada de por qué este diseño es apropiado para responder tu pregunta de investigación]}.

Según \cite{autor_diseños}, este tipo de diseño permite \textit{[ventajas específicas]}, lo cual es congruente con el objetivo de \textit{[tu objetivo]}.

% ----------------------------------------------------------------------------
% SECCIÓN 5.2: Población de Estudio
% ----------------------------------------------------------------------------
\section{Población de Estudio}
\label{sec:poblacion}

\subsection{Descripción de la Población}
\label{subsec:descripcion_poblacion}

La población de estudio estuvo conformada por \textit{[descripción clara de tu población objetivo]}.

\textbf{Características de la población:}

\begin{itemize}
    \item \textbf{Edad:} \textit{[rango etario]}
    \item \textbf{Sexo:} \textit{[hombres/mujeres/ambos]}
    \item \textbf{Ubicación:} \textit{[lugar específico]}
    \item \textbf{Contexto:} \textit{[hospital/escuela/comunidad/etc.]}
    \item \textbf{Período:} \textit{[fechas del estudio]}
\end{itemize}

\subsection{Criterios de Selección}
\label{subsec:criterios_seleccion}

\textbf{Criterios de Inclusión:}

Se incluyeron participantes que cumplieron con los siguientes criterios:

\begin{enumerate}
    \item \textit{[Criterio de inclusión 1]}
    \item \textit{[Criterio de inclusión 2]}
    \item \textit{[Criterio de inclusión 3]}
    \item \textit{[etc.]}
\end{enumerate}

\textbf{Criterios de Exclusión:}

Se excluyeron participantes que presentaron:

\begin{enumerate}
    \item \textit{[Criterio de exclusión 1]}
    \item \textit{[Criterio de exclusión 2]}
    \item \textit{[Criterio de exclusión 3]}
    \item \textit{[etc.]}
\end{enumerate}

\textbf{Criterios de Eliminación:}

Se eliminaron participantes cuando:

\begin{enumerate}
    \item \textit{[Criterio de eliminación 1, ej: no completaron el cuestionario]}
    \item \textit{[Criterio de eliminación 2, ej: retiraron su consentimiento]}
    \item \textit{[etc.]}
\end{enumerate}

\subsection{Tamaño de la Muestra}
\label{subsec:tamaño_muestra}

El tamaño de la muestra se calculó utilizando \textit{[fórmula/software]} con los siguientes parámetros:

\begin{itemize}
    \item \textbf{Nivel de confianza:} \textit{[95\% típicamente]}
    \item \textbf{Margen de error:} \textit{[5\% típicamente]}
    \item \textbf{Proporción esperada:} \textit{[valor o 50\% si desconoce]}
    \item \textbf{Tamaño de población:} \textit{[N total si es finita]}
\end{itemize}

\textbf{Cálculo del tamaño muestral:}

\begin{equation}
n = \frac{Z^2 \cdot p \cdot q}{d^2}
\label{eq:tamano_muestra}
\end{equation}

Donde:
\begin{itemize}
    \item $n$ = tamaño de la muestra
    \item $Z$ = nivel de confianza (1.96 para 95\%)
    \item $p$ = proporción esperada
    \item $q$ = 1 - p
    \item $d$ = precisión/margen de error
\end{itemize}

\textbf{Resultado:}

El tamaño muestral calculado fue de \textit{n = [número]} participantes. Se incluyeron finalmente \textit{[número real]} participantes.

\subsection{Tipo de Muestreo}
\label{subsec:tipo_muestreo}

Se utilizó un muestreo \textit{[probabilístico/no probabilístico]} de tipo \textit{[aleatorio simple/estratificado/por conglomerados/por conveniencia/intencional/etc.]}.

\textbf{Procedimiento de muestreo:}

\textit{[Describe CÓMO se seleccionaron los participantes]}

% ----------------------------------------------------------------------------
% SECCIÓN 5.3: Definición Operacional de Variables
% ----------------------------------------------------------------------------
\section{Definición Operacional de Variables}
\label{sec:variables}

% INSTRUCCIONES:
% - Esta sección es OBLIGATORIA y CRÍTICA
% - Presenta TODAS las variables de tu estudio
% - Usa una TABLA con formato APA
% - Incluye: definición semántica, operacional, dimensiones, escala de medición

\subsection{Variables del Estudio}
\label{subsec:variables_estudio}

Las variables consideradas en este estudio fueron:

\textbf{Variable Dependiente:}
\begin{itemize}
    \item \textit{[Nombre de la variable dependiente]}
\end{itemize}

\textbf{Variable(s) Independiente(s):}
\begin{itemize}
    \item \textit{[Variable independiente 1]}
    \item \textit{[Variable independiente 2]}
    \item \textit{[etc.]}
\end{itemize}

\textbf{Variables Universales (sociodemográficas):}
\begin{itemize}
    \item \textit{[Edad, sexo, escolaridad, etc.]}
\end{itemize}

\textbf{Variables Confusoras/Intervinientes (si aplica):}
\begin{itemize}
    \item \textit{[Variables que pueden afectar la relación estudiada]}
\end{itemize}

\subsection{Tabla de Operacionalización de Variables}
\label{subsec:operacionalizacion}

La \Cref{tab:operacionalizacion} presenta la definición operacional de todas las variables del estudio.

\begin{longtable}{@{}p{2.5cm}p{3cm}p{3cm}p{2cm}p{2cm}@{}}
    \caption{Operacionalización de variables del estudio} \label{tab:operacionalizacion} \\
    \toprule
    \textbf{Variable} & \textbf{Definición Semántica} & \textbf{Definición Operacional} & \textbf{Dimensiones} & \textbf{Escala de Medición} \\
    \midrule
    \endfirsthead
    
    \multicolumn{5}{c}{{\tablename\ \thetable{} -- Continuación}} \\
    \toprule
    \textbf{Variable} & \textbf{Definición Semántica} & \textbf{Definición Operacional} & \textbf{Dimensiones} & \textbf{Escala de Medición} \\
    \midrule
    \endhead
    
    \midrule
    \multicolumn{5}{r}{{Continúa en la siguiente página}} \\
    \endfoot
    
    \bottomrule
    \endlastfoot
    
    % VARIABLE DEPENDIENTE
    \textit{[Nombre VD]} & 
    \textit{[Definición conceptual según autor \cite{ref}]} & 
    \textit{[Cómo se medirá, qué instrumento, qué pregunta]} & 
    \textit{[Dimensiones o categorías si aplica]} & 
    \textit{[Nominal/ Ordinal/ Intervalo/ Razón]} \\
    \midrule
    
    % VARIABLES INDEPENDIENTES
    \textit{[Nombre VI1]} & 
    \textit{[Definición]} & 
    \textit{[Operacionalización]} & 
    \textit{[Dimensiones]} & 
    \textit{[Escala]} \\
    \midrule
    
    \textit{[Nombre VI2]} & 
    \textit{[Definición]} & 
    \textit{[Operacionalización]} & 
    \textit{[Dimensiones]} & 
    \textit{[Escala]} \\
    \midrule
    
    % VARIABLES UNIVERSALES
    Edad & 
    Tiempo transcurrido desde el nacimiento & 
    Años cumplidos reportados por el participante & 
    Años & 
    Razón (continua) \\
    \midrule
    
    Sexo & 
    Condición biológica & 
    Autorreporte: Hombre/Mujer & 
    Hombre, Mujer & 
    Nominal (dicotómica) \\
    \midrule
    
    \textit{[Continúa con todas tus variables]} & 
    ... & 
    ... & 
    ... & 
    ... \\
    
\end{longtable}

% ----------------------------------------------------------------------------
% SECCIÓN 5.4: Técnicas y Procedimiento
% ----------------------------------------------------------------------------
\section{Técnicas y Procedimiento}
\label{sec:tecnicas_procedimiento}

\subsection{Instrumentos de Recolección de Datos}
\label{subsec:instrumentos}

Para la recolección de datos se utilizaron los siguientes instrumentos:

\textbf{Instrumento 1: \textit{[Nombre del instrumento]}}

\begin{itemize}
    \item \textbf{Autor:} \textit{[Autor original]} \cite{referencia}
    \item \textbf{Propósito:} \textit{[Qué mide]}
    \item \textbf{Estructura:} \textit{[Número de ítems, dimensiones]}
    \item \textbf{Escala de respuesta:} \textit{[Likert/Dicotómica/etc.]}
    \item \textbf{Puntuación:} \textit{[Cómo se puntúa, rangos]}
    \item \textbf{Validez y confiabilidad:} \textit{[Reportar del estudio original]}
\end{itemize}

\textbf{Instrumento 2: \textit{[Si aplica]}}

[Repite la estructura]

\subsection{Validación de Instrumentos}
\label{subsec:validacion_instrumentos}

% SOLO SI APLICA: Si validaste o adaptaste instrumentos

Los instrumentos fueron sometidos a procesos de validación que incluyeron:

\textbf{1. Validación de Contenido (Juicio de Expertos):}

Se solicitó la evaluación de \textit{[número]} expertos en \textit{[área]}, quienes valoraron:
\begin{itemize}
    \item Claridad de los ítems
    \item Pertinencia con las variables
    \item Suficiencia del instrumento
\end{itemize}

Se calculó el Índice de Validez de Contenido (IVC) obteniendo un valor de \textit{[valor]}, considerado \textit{[aceptable/excelente según criterio]}.

\textbf{2. Análisis de Confiabilidad (Consistencia Interna):}

Se aplicó el coeficiente Alfa de Cronbach para evaluar la consistencia interna, obteniendo:

\begin{itemize}
    \item \textbf{Instrumento global:} $\alpha$ = \textit{[valor]}
    \item \textbf{Dimensión 1:} $\alpha$ = \textit{[valor]}
    \item \textbf{Dimensión 2:} $\alpha$ = \textit{[valor]}
\end{itemize}

Valores de $\alpha \geq 0.70$ se consideraron aceptables \cite{referencia_metodologica}.

\textbf{3. Validación de Constructo (si muestra amplia):}

Se realizó Análisis Factorial Exploratorio (AFE) con los siguientes resultados:

\begin{itemize}
    \item \textbf{KMO:} \textit{[valor]} (\textit{[interpretación]})
    \item \textbf{Prueba de Bartlett:} $\chi^2 = $ \textit{[valor]}, $p < 0.001$
    \item \textbf{Varianza explicada:} \textit{[\%]}
    \item \textbf{Factores extraídos:} \textit{[número]}
\end{itemize}

\subsection{Procedimiento de Recolección de Datos}
\label{subsec:procedimiento_recoleccion}

% INSTRUCCIONES:
% - Describe PASO A PASO el procedimiento
% - Sé ESPECÍFICO: cuándo, dónde, cómo, quién
% - Redacta en PASADO
% - Incluye reclutamiento, medición, técnicas empleadas

El procedimiento de recolección de datos se desarrolló de la siguiente manera:

\textbf{Fase 1: Reclutamiento de Participantes (\textit{[fechas]})}

\begin{enumerate}
    \item Se \textit{[verbo en pasado: contactó/visitó/convocó]} a \textit{[población]} en \textit{[lugar]}
    \item Se explicó el propósito del estudio y se entregó la carta de consentimiento informado
    \item Se verificaron los criterios de inclusión/exclusión
    \item Se obtuvo el consentimiento informado por escrito
\end{enumerate}

\textbf{Fase 2: Aplicación de Instrumentos (\textit{[fechas]})}

\begin{enumerate}
    \item Se aplicaron los instrumentos en \textit{[modalidad: presencial/virtual]}
    \item El tiempo promedio de aplicación fue de \textit{[minutos]} minutos
    \item Los instrumentos fueron administrados por \textit{[quién: investigador/personal capacitado]}
    \item Se realizó en \textit{[lugar específico]} bajo condiciones de \textit{[privacidad/comodidad/etc.]}
\end{enumerate}

\textbf{Fase 3: Técnicas Complementarias (\textit{[si aplica]})}

[Si usaste técnicas de laboratorio, clínicas, bioquímicas, etc., descríbelas aquí]

\begin{itemize}
    \item \textbf{Técnica 1:} \textit{[Descripción detallada del procedimiento]}
    \item \textbf{Equipo utilizado:} \textit{[Especificaciones]}
    \item \textbf{Personal:} \textit{[Quién realizó, capacitación]}
\end{itemize}

\textbf{Fase 4: Control de Calidad}

\begin{itemize}
    \item Se verificó la completitud de los instrumentos
    \item Se codificaron las respuestas según \textit{[criterio]}
    \item Se capturaron los datos en \textit{[software]} por duplicado para verificar exactitud
\end{itemize}

\subsection{Asignación de Grupos (Solo para Estudios de Intervención)}
\label{subsec:asignacion_grupos}

% SOLO SI APLICA

Los participantes fueron asignados a los grupos de estudio mediante \textit{[método de aleatorización/asignación]}.

\textbf{Grupos conformados:}

\begin{itemize}
    \item \textbf{Grupo Experimental:} \textit{n = [número]} participantes que recibieron \textit{[intervención]}
    \item \textbf{Grupo Control:} \textit{n = [número]} participantes que recibieron \textit{[placebo/tratamiento estándar/nada]}
\end{itemize}

\textbf{Cegamiento:}

El estudio fue \textit{[simple ciego/doble ciego/abierto]}, donde \textit{[quien fue cegado: participantes/investigadores/ambos]}.

% ----------------------------------------------------------------------------
% SECCIÓN 5.5: Análisis de la Información
% ----------------------------------------------------------------------------
\section{Análisis de la Información}
\label{sec:analisis_informacion}

\subsection{Software Estadístico}
\label{subsec:software}

El análisis de datos se realizó utilizando el software \textit{[SPSS/R/Stata/etc.]} versión \textit{[número]}.

\subsection{Análisis Univariado}
\label{subsec:analisis_univariado}

Se realizó análisis descriptivo de todas las variables:

\textbf{Variables cuantitativas:}
\begin{itemize}
    \item Medidas de tendencia central: media, mediana
    \item Medidas de dispersión: desviación estándar, rango intercuartílico
    \item Se evaluó normalidad mediante \textit{[Kolmogorov-Smirnov/Shapiro-Wilk]}
\end{itemize}

\textbf{Variables cualitativas:}
\begin{itemize}
    \item Frecuencias absolutas y relativas
    \item Porcentajes e intervalos de confianza del 95\%
\end{itemize}

\subsection{Análisis Bivariado}
\label{subsec:analisis_bivariado}

Para evaluar la relación entre variables se utilizaron:

\textbf{Variables cuantitativas vs cuantitativas:}
\begin{itemize}
    \item Coeficiente de correlación de Pearson (si distribución normal)
    \item Coeficiente de Spearman (si distribución no normal)
\end{itemize}

\textbf{Variables cualitativas vs cualitativas:}
\begin{itemize}
    \item Chi-cuadrada ($\chi^2$) o Prueba exacta de Fisher
    \item Cálculo de Odds Ratio (OR) e intervalos de confianza
\end{itemize}

\textbf{Variables cuantitativas vs cualitativas:}
\begin{itemize}
    \item Prueba t de Student (si 2 grupos, distribución normal)
    \item Prueba U de Mann-Whitney (si 2 grupos, no normal)
    \item ANOVA (si >2 grupos, normal) o Kruskal-Wallis (no normal)
\end{itemize}

Se consideró significancia estadística cuando $p < 0.05$.

\subsection{Análisis Multivariado}
\label{subsec:analisis_multivariado}

Para controlar variables confusoras y evaluar predictores se realizó:

\textbf{[Según tu estudio, puede ser:]}

\begin{itemize}
    \item \textbf{Regresión lineal múltiple:} Para predecir \textit{[variable dependiente continua]}
    \item \textbf{Regresión logística:} Para predecir \textit{[variable dependiente dicotómica]}
    \item \textbf{Análisis de covarianza (ANCOVA):} Para controlar \textit{[covariables]}
    \item \textbf{Análisis factorial:} Para identificar estructura subyacente de variables
\end{itemize}

Se reportaron coeficientes, intervalos de confianza del 95\% y valores de $p$.

% ----------------------------------------------------------------------------
% SECCIÓN 5.6: Consideraciones Éticas
% ----------------------------------------------------------------------------
\section{Consideraciones Éticas}
\label{sec:consideraciones_eticas}

% INSTRUCCIONES:
% - Esta sección es OBLIGATORIA
% - Especifica lineamientos internacionales, nacionales, locales
% - La carta de consentimiento va en ANEXOS

\subsection{Principios Éticos}
\label{subsec:principios_eticos}

Esta investigación se desarrolló respetando los principios éticos establecidos en:

\begin{itemize}
    \item \textbf{Declaración de Helsinki} \cite{helsinki}: Principios éticos para investigación médica en seres humanos
    \item \textbf{Informe Belmont} \cite{belmont}: Principios de respeto a las personas, beneficencia y justicia
    \item \textbf{Reglamento de la Ley General de Salud en Materia de Investigación para la Salud (México)}: Artículos \textit{[específicos que apliquen]}
    \item \textbf{Normas institucionales de la UACH}: \textit{[Si aplica]}
\end{itemize}

\subsection{Aspectos Éticos Específicos}
\label{subsec:aspectos_eticos_especificos}

\textbf{Consentimiento Informado:}

\begin{itemize}
    \item Se obtuvo consentimiento informado por escrito de todos los participantes
    \item El documento explicó claramente: objetivo, procedimientos, riesgos, beneficios, confidencialidad, voluntariedad
    \item Los participantes tuvieron oportunidad de hacer preguntas
    \item Se enfatizó el derecho a retirarse en cualquier momento sin consecuencias
    \item [La carta de consentimiento se presenta en el Anexo \ref{anexo:consentimiento}]
\end{itemize}

\textbf{Confidencialidad y Privacidad:}

\begin{itemize}
    \item Los datos personales fueron codificados para proteger la identidad
    \item La información se almacenó en \textit{[medio seguro: computadora con contraseña/archivo bajo llave]}
    \item Solo el equipo de investigación tuvo acceso a los datos
    \item Los resultados se reportan de forma agregada, sin identificadores personales
\end{itemize}

\textbf{Riesgos y Beneficios:}

\begin{itemize}
    \item \textbf{Riesgos:} \textit{[Especifica los riesgos, pueden ser mínimos]}
    
    Ejemplo: \textit{La participación en este estudio representó riesgos mínimos, limitados a molestia o cansancio al responder los cuestionarios.}
    
    \item \textbf{Beneficios:} \textit{[Especifica beneficios directos o indirectos]}
    
    Ejemplo: \textit{Los participantes recibieron retroalimentación sobre sus resultados y recomendaciones personalizadas.}
\end{itemize}

\textbf{Voluntariedad:}

\begin{itemize}
    \item La participación fue completamente voluntaria
    \item No hubo coerción ni incentivos indebidos
    \item Los participantes pudieron retirarse en cualquier momento
\end{itemize}

\subsection{Aprobación por Comité de Ética}
\label{subsec:aprobacion_comite}

\textbf{[Si aplica:]}

Este protocolo de investigación fue sometido a evaluación y aprobado por el Comité de Ética en Investigación de \textit{[nombre de la institución]} con número de registro \textit{[número]} en fecha \textit{[fecha]}.

\textbf{[Si NO aplica:]}

Este proyecto no requirió aprobación de comité de ética porque \textit{[razón: investigación con bajo riesgo, uso de datos secundarios anónimos, etc.]}, sin embargo, se siguieron todos los principios éticos mencionados.

% ============================================================================
% NOTAS FINALES:
% ============================================================================
%
% RECUERDA:
% 1. TODO en TIEMPO PASADO (ya se hizo)
% 2. Tabla de operacionalización es OBLIGATORIA
% 3. Procedimiento debe ser DETALLADO y REPLICABLE
% 4. Consideraciones éticas son OBLIGATORIAS
% 5. Consentimiento informado va en ANEXOS
% 6. Análisis de datos por NIVELES (univariado, bivariado, multivariado)
%
% ============================================================================

